\begin{table}[t]
\centering
\small
\caption{Cross-benchmark evaluation with GPT-4o. ASB results are averaged across 3 runs (565 cases). InjecAgent (60 attack-only cases) and AgentDojo (510 cases) are evaluated via adapter conversion.}
\label{tab:cross_benchmark}
\begin{tabular}{l cc cc}
\toprule
 & \multicolumn{2}{c}{\textbf{InjecAgent}} & \multicolumn{2}{c}{\textbf{AgentDojo}} \\
\cmidrule(lr){2-3} \cmidrule(lr){4-5}
\textbf{Defense} & ASR & BSR & ASR & BSR \\
\midrule
D0 (Baseline)    & 0.000 & ---   & 0.000 & 1.000 \\
D1 (Spotlighting)& 0.000 & ---   & 0.000 & 1.000 \\
D5 (Sandwich)    & 0.000 & ---   & 0.000 & 1.000 \\
D10 (\civ{})     & 0.000 & ---   & 0.000 & 1.000 \\
\bottomrule
\end{tabular}
\begin{flushleft}
\scriptsize All defenses achieve 0\% ASR on both external benchmarks, indicating that the adapter-converted attack scenarios do not trigger forbidden actions under ASB's rule-based judge.
The zero ASR for the D0 baseline suggests a mismatch between the external benchmarks' attack success criteria and ASB's evaluation protocol, rather than genuine defense effectiveness.
This highlights the need for benchmark-native evaluation when comparing across frameworks.
\end{flushleft}
\end{table}
